\documentclass[10pt,a4paper]{amsart}
\usepackage[margin=19mm]{geometry}
\usepackage{amsmath,amssymb,mathtools}
\usepackage[scr=rsfs,cal=euler]{mathalfa}
\usepackage{xfrac}
\usepackage[unicode,hidelinks,bookmarks=false]{hyperref}
\newcommand{\ie}{i.e.\ }
\newcommand{\cf}{cf.\ }
\newcommand{\resp}{respectively\ }
\newcommand{\Subsection}{\subsection}
\providecommand{\nobreakdash}{\nobreak\hbox{-}}
\setlength{\emergencystretch}{2em}
\multlinegap0pt
\usepackage{amssymb}
\usepackage{mathtools}
\usepackage{bm}
\usepackage{float}
\usepackage{tikz}
\usepackage{csquotes}
\usepackage{caption}
\newtheorem{lemma}{Lemma}
\newtheorem{proposition}{Proposition}
\newtheorem{definition}{Definition}
\newtheorem{theorem}{Theorem}
\newtheorem{claim}{Claim}
     \newcommand{\bT}{\mathbb{T}}
  \newcommand{\calG}{\mathcal{G}}
  \newcommand{\calN}{\mathcal{N}}
    \newcommand{\calY}{\mathcal{Y}}
  \newcommand{\calZ}{\mathcal{Z}}
\newcommand{\diam}{\mathrm{diam}}
\title{Bicombing the mapping class group and Teichm\"uller space via stable cubical intervals}
\author{Matthew Gentry Durham}
\date{Source: arXiv:2512.24136v1 (CC BY 4.0)}
\begin{document}
\maketitle

\noindent\emph{Source and licence.} Bicombing the mapping class group and Teichm\"uller space via stable cubical intervals. Version arXiv:2512.24136v1 (CC BY 4.0), distributed by arXiv under the Creative Commons Attribution 4.0 International licence, http://creativecommons.org/licenses/by/4.0/, as recorded in the arXiv licence metadata for this exact version and shown on the abstract page https://arxiv.org/abs/2512.24136v1. Full licence notice: \texttt{licenses/arxiv-2512.24136-CC-BY-4.0.txt}.\par\smallskip

\noindent\textit{Editorial scope.} Counted unit: the article's proposition cluster (source0.tex lines 2632--2640). The article's complete original proof of it is delivered with it (source0.tex lines 2642--2695, one \texttt{proof} environment of 54 lines). No mathematics is written, repaired or re-derived: the statement and the proof are the article's own lines, in the article's own notation, and no retained line is substituted or re-flowed.  Retained uncounted as the source-local context the proof needs: lines 1943--1949; lines 1966--1968; lines 1989--1993; lines 2033--2042; lines 2096--2103; lines 2113--2138; lines 2150--2160; lines 2239--2243; lines 2246--2251.  One declared adaptation: exactly 2 ancillary strings inside the retained spans are omitted (image-inclusion commands and, where the source repeats a label, the repeated label definitions) and no other line differs from the source; the figure environments, with their placement, captions and labels, are kept, and the package records the omitted strings in fidelity receipts bound to the affected bodies.  No retained line contains a citation, so the wrapper carries no bibliography and no bibliography entry.  The retained lines contain the article's own diagram code, which the wrapper typesets with the standard tikz packages; no image file is delivered and the package declares no asset dependency.  Editorial formatting only, in addition: an AMS \texttt{amsart} preamble that restates the article's own declarations and macros used by the retained lines, namely the environments \texttt{lemma}, \texttt{proposition}, \texttt{definition}, \texttt{theorem}, \texttt{claim} on separate counters, together with the standard packages those lines need.  The retained environments print in this document as Lemma 1, Proposition 1, Definition 1, Theorem 1, Claim 1 in order of appearance, whereas the article prints the same items with the numbering of its own full document.  The complete native source of the article as distributed in the arXiv e-print for this exact version is bundled unchanged as \texttt{sources/191856/source0.tex}.
\begin{lemma}[Fellow-traveling, $\delta$-hyperbolic]\label{lem:Morse close}
Let $\calZ$ be $\delta$-hyperbolic and geodesic, and $\gamma$ a geodesic between $a,b \in \calZ$.  For every $\epsilon>0$ there exists $\epsilon'=O(\epsilon)>0$ so that the following holds:

\begin{itemize}
\item Suppose $\gamma'$ is a geodesic connecting points $x,y \in \calN_{\epsilon}(\gamma)$, and that $x',y'$ are closest points to $x,y$ on $\gamma$, respectively.  Then any geodesic $[x,y]$ satisfies $d_{\calZ}^{Haus}([x,y], [x',y']) < \epsilon'$.
\end{itemize}
\end{lemma}

\begin{lemma} \label{lem:cluster arrangement}
If $E > 5\epsilon$ and $C \neq C'$ are distinct $E$-clusters, then $s(C) \cap s(C') = \emptyset$.
\end{lemma}

\begin{proposition}\label{prop:cluster sep graph}
Let $\calZ$ be $\delta$-hyperbolic, $\lambda(a,b)$ a geodesic connecting $a,b \in \calZ$, and $\calY \subset \calZ$ finite with $\calY \subset \calN_{\epsilon}(\lambda(a,b))$.  Also let $\epsilon' = \epsilon'(\epsilon, \delta)>0$ as in Lemma \ref{lem:Morse close}, with also $\epsilon' > \epsilon$.

There exists $E = E(\epsilon, \delta)>0$ so that the $E$-cluster separation graph $\calG_{E,\epsilon'}(a,b;\calY)$ is an interval whose order coincides with the order of cluster shadows along $\lambda(a,b)$, with the endpoint clusters containing $a,b$, respectively. 
\end{proposition}

\begin{lemma}\label{lem:cluster close}
Let $C_1, C_2$ be distinct clusters along $\lambda(a,b)$.  Suppose that $p_1 \in C_1$ and $p_2 \in C_2$ are closest points, while $x_1 \in s(C_1)$ and $x_2 \in s(C_2)$ are closest points.  Then $d_{\calZ}(x_i, p_i) \leq 7\epsilon$ for $i=1,2$.

As a consequence, the following hold:
\begin{enumerate}
\item For distinct clusters $C_1,C_2$, we have $|d(C_1,C_2) - d(s(C_1),s(C_2))| < 14\epsilon.$
\item For any $p \in T_e$, we have that $d_{\calZ}(p, \calY \cup \{a,b\}) \leq d_{\calZ}(p, \partial T_e) + O(\epsilon).$
\item For any components $V,W$ of $T = T_e \cup T_c$, we have $\diam_{\calZ}(s(V)\cap s(W)) < O(\epsilon)$.
\end{enumerate}
\end{lemma}

\begin{definition}[Order-preserving maps]\label{defn:order preserving}
Let $T_1, T_2$ be two intervals with the endpoints of $T_i$ labeled by $a_i,b_i$.  Suppose that $T_i = T_{i,1} \sqcup T_{i,2}$ are subdivisions of $T_i$ into finite collections of subintervals with $\# T_{1,1} = \#T_{2,1}$ (and hence $\#T_{1,2} = \#T_{2,2}$).

\begin{itemize}
\item We say an isometry $i_{E,E'}:E \to E'$ between components $E \subset T_{1,1}$ and $E' \subset T_{2,1}$ is \emph{order preserving} if $i_{E,E'}$ sends the endpoint of $E$ closest to $a_1$ to the endpoint of $E'$ closest to $a_2$.
\item We say a bijection $\alpha:\pi_0(T_{1,1}) \to \pi_0(T_{2,1})$ is \emph{order-preserving} if it coincides with the bijection obtained by recording the orders of appearance of the components of $T_{i,1}$ along $T_i$ from $a_i$ to $b_i$.
\end{itemize}
\end{definition}

\begin{definition}[Stable decomposition] \label{defn:stable decomp}
Let $\calZ$ be $\delta$-hyperbolic and geodesic.  
\begin{enumerate}
\item Given $N>0$, two $\epsilon$-setups $(a,b; \calY)$ and $(a',b'; \calY')$ are $(N,\epsilon)$-\emph{admissible} if
\begin{enumerate}
\item $d_{\calZ}(a,a'), d_{\calZ}(b,b') \leq \epsilon$,
\item $\calY, \calY' \subset \calN_{\epsilon/2}(\lambda(a,b)) \cap \calN_{\epsilon/2}(\lambda(a',b'))$,
\item $|\calY \triangle \calY'| < N$.
\end{enumerate}
\item Given an $\epsilon$-setup $(a,b;\calY)$, an \emph{edge decomposition} of its stable interval $T = T_e \cup T_c$ is a collection of subintervals $T_s \subset T_e$.
\item Given $L>0$, $\calY_0 \subset \calY \cap \calY'$, and two $(N,\epsilon)$-admissible setups $(a,b;\calY)$ and $(a',b'; \calY')$, we say that two edge decompositions $T_s \subset T_e$ and $T'_s \subset T'_e$ are $\calY_0$-\emph{stably} $L$-\emph{compatible} if
\begin{enumerate}
\item Each component of $T_s$ and $T'_s$ has positive integer length with endpoints at vertices of $T_e$ and $T'_e$, respectively. \label{item:integer length}
\item There is an order-preserving bijection $\alpha:\pi_0(T_s) \to \pi_0(T'_s)$ between the sets of stable components. \label{item:stable bijection}
\item For each \emph{stable pair} $(E, E')$ identified by $\alpha$, there exists an order-preserving isometry $i_{E,E'}:E \to E'$. \label{item:stable pairs}
\item For all but at most $L$ pairs of stable components $(E,E')$, we have $\phi(E) = \phi(E')$ are identical with $\phi(x) = \phi'(i_{E,E'}(x))$ for all $x \in E$. \label{item:identical pairs}
\item For the (at most) $L$-many remaining stable pairs $(E,E')$, we have $d_{\calZ}(\phi(x), \phi'(i_{E,E'}(x)))< L$ for all $x \in E$. \label{item:close pairs}
\item The complements $T_{e} - T_s$ and $T'_{e} - T'_s$ consist of at most $L$-many \emph{unstable components} of diameter at most $L$.\label{item:unstable components}
\item The induced order-preserving bijection $\beta:\pi_0(T-T_s) \to \pi_0(T' - T'_s)$ satisfies: \label{item:adjacency}
\begin{enumerate}
\item (Endpoints) Let $D_a$ denote the component of $T-T_s$ containing $\mu(C_a)$, and define $D_b,D_{a'},D_{b'}$ similarly.  Then $\beta(D_a) = D_{a'}$ and $\beta(D_b) = D_{b'}$. \label{item:endpoint condition}
\item (Identifying clusters) For any $y \in \calY_0$, let $D_y, D'_y$ denote the components of $T-T_s$ and $T' - T'_s$ containing $\mu(C_y), \mu(C'_y)$, respectively.  Then $\beta(D_y) = D'_y$. \label{item:cluster identify}
\end{enumerate}
\end{enumerate}
\end{enumerate}
\end{definition}

\begin{theorem}[Stable intervals]\label{thm:stable intervals}
Let $\calZ$ be $\delta$-hyperbolic and geodesic.  For any $\epsilon, N>0$ and positive integers $r_1,r_2>0$, there exist $L_1 = L_1(\delta, \epsilon, N)>0$ and $L_2 = L_2(\delta, \epsilon, N,r_1,r_2)>0$ so that the following holds.  Suppose $(a,b;\calY)$ and $(a',b';\calY')$ are $(N,\epsilon)$-admissible $\epsilon$-setups and let $T = T_e \cup T_c$ and $T' = T'_e \cup T'_c$ denote their stable intervals.  Then
\begin{enumerate}
\item There exist $(\calY \cap \calY')$-stable $L_1$-compatible decompositions $T_s \subset T_e$ and $T'_s \subset T'_e$, and
\item There exist $(\calY \cap \calY')$-stable $L_2$-compatible edge decompositions $\bT_{s} \subset \bT_{e}$ and $\bT'_{s} \subset \bT'_{e}$ of the $(r_1,r_2)$-thickenings of $T, T'$ along $T_c,T'_c$.
\begin{itemize}
\item Moreover, we have $\bT_s \subset T_s$ and $\bT_s \subset T'_s$.
\end{itemize}
\end{enumerate}

\end{theorem}

\begin{figure}[b]
    \centering
    
    \caption{Step (1) of the proof of Theorem \ref{thm:stable intervals}, split cluster case:  In the first figure, the two edge components of the two stable intervals are shown.  The point $z$ forms an individual cluster and therefore must be inserted as a vertex in the graph $\calG$, between clusters $C_1, C_2$.  Every cluster to the left of $C_1$ is unchanged, as are all of the connecting segments (in pink), and same for those to the right of $C_2$.  The second figure shows the stable decomposition of the stable intervals.  Here, we need to include the shadow of the cluster $z$ on $\lambda(C_1,C_2)$ as an unstable component.  This cuts $\lambda(C_1,C_2)$ into two segments, which are paired with $\lambda(C_1,z)$ and $\lambda(z,C_2)$, though these need to be slightly truncated so guarantee they have the same lengths; the unstable segments accomplishing this are red.  The third figure shows the stable decomposition for their $(r_1,r_2)$-thickenings, which are subsets of the edge components $T_e \subset T$ and $T'_e \subset T$, are indicated in brown.  Note that this can result in collapsing an entire edge component if its diameter is less than $r_2-2r_1$.  These brown segments refine the stable decompositions $T_s \subset T_e$ and $T'_s \subset T'_e$.}    \label{fig:split cluster}
\end{figure}

\begin{figure}[!htb]
    \centering
    
    \caption{Step (1) of the proof of Theorem \ref{thm:stable intervals}, affected cluster case.  Two phenomena are shown.  In the top figure, there are three affected clusters, $P,Q,R$, for $(a,b;\calY)$ which combine with $z$ into a new cluster $C'$ for $(a,b;\calY \cup \{z\})$.  By definition, $P,Q,R$ must be $O(\epsilon)$ apart, and actually are adjacent in $\calG$.  The clusters $C_1,C_2$ appear in both setups, and the subgraph of $\calG$ to the left of $C_1$ coincides with the same subgraph for $\calG'$, and similarly for $C_2$.  So the corresponding (pink) components form (stable) components of $T_s \subset T_e$ and $T'_s \subset T'_e$.  The components of $T_e$ connecting $P,Q,R$ are declared unstable and, in this example, the components $\lambda(C_1,P)$ and $\lambda(R,C_2)$ of $T_e$ coincide with the components $\lambda(C_1,C')$ and $\lambda(C',C_2)$ of $T'_e$, and hence are stable.  The second figure shows the stable decomposition for the $(r_1,r_2)$-thickenings, which are just the thickenings of the components of $T_s = T'_s$, and hence are also exactly the same.  In the second example, the same discussion applies to the outside parts.  For the middle of the graph, $z$ combines with the cluster $A$, which results in pairs of segments $\lambda(C_1,A)$ and $\lambda(C_1, C')$, and $\lambda(A,C_2)$ and $\lambda(C',C_2)$, which have $O(\epsilon)$-close endpoints and are thus $O(\epsilon)$-Hausdorff close, with the same length up to error $O(\epsilon)$.  Small surgeries provide the required stable components for the stable intervals $T, T'$, shown in the third figure.  The fourth figure shows the stable decomposition for the $(R,r)$-thickenings of this example.}
    \label{fig:affected cluster}
\end{figure}

\begin{proposition} \label{prop:add cluster}
Let $\calZ$ be a geodesic $\delta$-hyperbolic space, $\epsilon>0$, and $r_1,r_2>0$ be positive integers.  Let $(a,b;\calY)$ be an $\epsilon$-setup and further fix a point $z \in \calN_{\epsilon/2}(\lambda(a,b))$.  Then
\begin{enumerate}
\item $(a,b;\calY \cup \{z\})$ is an $\epsilon$-setup;
\item The two setups $(a,b;\calY)$ and $(a,b;\calY \cup \{z\})$ are $(1,\epsilon)$-admissible;
\item There exists $L_1=L_1(\delta, \epsilon)>0$ so that the stable intervals $T$ and $T'$ for the setups $(a,b;\calY)$ and $(a,b;\calY \cup \{z\})$ admit $\calY$-stable $L_1$-compatible decompositions;
\item There exists $L_2 = L_2(\delta, \epsilon, r_1,r_2)>0$ so that the $(r_1,r_2)$-thickenings of $T,T'$ admit $\calY$-stable $L_2$-compatible decompositions.
\end{enumerate}
\end{proposition}

\begin{proof}
Since the points $a,b$ have not changed, the first two conclusions are automatic from the assumptions.

Choosing $E \gg \epsilon$, Proposition \ref{prop:cluster sep graph} implies that both $\calG = \calG_E(a,b;\calY)$ and $\calG' = \calG_E(a,b;\calY \cup \{z\})$ are intervals.  There are two main cases: (a) $d(z,C)>E$ for all clusters $C$, or (b) otherwise.

We will deal with these cases separately, though the arguments are similar.  Moreover, for each case, we will want to deal both with the original decompositions of the stable tree $T = T_e \cup T_c$, as well as its $(r_1,r_2)$-thickening $T = \bT_{e} \cup \bT_{c}$.  The compatible decompositions for the latter will be refined for the decomposition for the former, with possibly a slight increase in the error depending on $r_1,r_2$.

\medskip

\underline{\textbf{Case (a), $z$ forms a split cluster}}: In this case, $z$ forms its own cluster and, in particular, does not affect membership of any other cluster.  Moreover, it cannot be an endpoint cluster of $\calG'$, as those necessarily contain $a$ and $b$.  Hence there exist two clusters $C_1,C_2$ which are adjacent in $\calG$ but are separated by $z$ in $\calG'$, say with $C_1$ on the side of $a$ and $C_2$ on the side of $b$.  In this sense, we think of $z$ as splitting the clusters $C_1,C_2$.

Lemma \ref{lem:cluster arrangement} implies that $z$ separates two clusters $C,C'$ for $(a,b;\calY)$ if and only if $s(z)$ separates $s(C), s(C')$ along $\lambda(a,b)$.  It follows then that the subgraph of $\calG$ connecting the cluster containing $a$ to $C_1$ is identical to the subgraph of $\calG_E(a,b;\calY \cup \{z\})$, and similarly for the subgraphs connecting $C_2$ to the cluster containing $b$.  This says that for any pair $C,C'$ on the $a$-side of $z$ in $\calG'$, $\lambda_0(C,C')$ is an edge component of both $T$ and $T'$, and similarly for such pairs  on the $b$-side of $z$.  Thus we can include each such component of $E \subset T_e$ and $E \subset T'_e$ in the respective stable decomposition $T_s \subset T_e$ and $T'_s \subset T'_e$, where the identity map is the map described in parts \eqref{item:stable pairs} and \eqref{item:identical pairs} of Definition \ref{defn:stable decomp}.  By construction, each of these pieces has positive integer length, thus verifying property \eqref{item:integer length}.

To construct the rest of the stable decompositions, namely stable pieces of type \eqref{item:close pairs} in Definition \ref{defn:stable decomp}, we need to compare the component $\lambda_0(C_1, C_2) \subset T_e$ to the two components $\lambda_0(C_1,z) \cup \lambda_0(z, C_2) \subset T'_e$.  First, Lemma \ref{lem:cluster close} says that the endpoints of $\lambda(C_1,C_2)$ and $\lambda(C_1,z)$ on $C_1$ are $14\epsilon$-close, and hence the endpoints of $\lambda_0(C_1,C_2)$ and $\lambda_0(C_1,z)$ are within $14\epsilon+2$, and similarly for $C_2$.  It follows from $\delta$-hyperbolicity of $\calZ$ that $\lambda_0(C_1, C_2)$ and $ \lambda_0(C_1,z) \cup \lambda_0(z, C_2)$ are $O(\epsilon)$-Hausdorff close.  Now, unless $C_1,C_2$ are $O(\epsilon)$-close in $\calZ$---in which case we can declare all three of these geodesics to be unstable---then the shadow $s_T(z)$ of $z$ on $T$ will be properly contained in $\lambda_0(C_1,C_2)$ and have $O(\epsilon)$-bounded diameter.    Hence we can decompose these geodesics into $O(\epsilon)$-close stable pieces and unstable pieces labeled by $z$, as described in property \eqref{item:unstable components}, while satisfying the length requirement of \eqref{item:integer length} and proximity requirements of \eqref{item:identical pairs} and \eqref{item:close pairs}; see Figure \ref{fig:split cluster} for a heuristic picture.  For the bijection $\alpha:\pi_0(T_s) \to \pi_0(T'_s)$, since $T_s,T'_s$ have the same number of components, we can take let $\alpha$ denote the natural order-preserving bijection, i.e. the one obtained by moving along the geodesics $[a,b]_{T}$ and $[a,b]_{T'}$.  By construction, this bijection identifies the components we want so that the related conditions \eqref{item:identical pairs} and \eqref{item:close pairs} are satisfied.

Finally, we need to confirm the properties in \eqref{item:adjacency} for the order-preserving bijection $\beta:\pi_0(T - T_s) \to \pi_0 (T'-T'_s)$ induced by $\alpha$.  Since $\beta$ is order-preserving (with the order being ``from $a$ to $b$'', as in Definition \ref{defn:order preserving}), \eqref{item:endpoint condition}, i.e. that the components of $T-T_s$ and $T' - T'_s$ associated to $a$ and $b$ are identified, is automatically satisfied.  For \eqref{item:cluster identify}, note that $\calY \cap (\calY \cup \{z\}) = \calY$, and, by construction again, if $y \in \calY$, then $\beta$ identifies the component $D_y$ of $T- T_s$ corresponding to $y$ with $D'_y$, the component of $T'-T'_s$ corresponding to $y$.  In this case, this is because all such components are exactly the same.   This proves part (1) of the statement.

For part (2) of the statement, we want to consider how the $(r_1,r_2)$-thickenings of $T,T'$ along $T_c,T'_c$ respectively affect the decompositions $T_s \subset T_e$ and $T'_s \subset T'_e$.  First, observe that $\bT_e \subset T_e$ and $\bT'_e \subset T'_e$ by definition.  In particular, consider a component $D \subset T_e$, where $D = \lambda_0(C,C')$ for adjacent clusters $C,C'$.  Then the $(r_1,r_2)$-thickening of $D$ either truncates $D$ to $D_{thick} = D - (\mu_{r_1}(C) \cup \mu_{r_1}(C'))$, or possibly removes $D$ entirely from $\bT_e$ if $d_T(\mu(C), \mu(C'))<r_2-2r_1$.  A similar statement holds for components of $T'_e$.

By the above, most components of $T_s$ and $T'_s$ are identical components of $T_e, T'_e$, and hence their thickenings are exactly the same by the above.  For the remaining components, namely $\lambda_0(C_1,z), \lambda_0(z,C_2)$ of $T'_e$ and $\lambda_0(C_1,C_2)$ of $T_e$, each is either truncated or removed as above.  The remaining two components of the stable decomposition $T_s$ correspond to $\lambda_0(C_1,z)$ and $\lambda_0(z,C_2)$, and so one can construct the corresponding components of $\bT_s$ by applying the same procedure to the $(r_1,r_2)$-thickening of $\lambda_0(C_1,z)$ and $\lambda_0(z,C_2)$.  Note that this can inflate the proximity constants, but only in terms of $r_1,r_2$.  The rest of the proof goes through unchanged; see Figure \ref{fig:split cluster}.




\medskip

\underline{\textbf{Case (b), $z$ affects other clusters}}: Suppose now instead that there exists some cluster $C$ for $(a,b;\calY)$ so that $d(z,C)<E$; we say $C$ is an  \emph{affected cluster}.

\begin{claim}\label{claim:bounded affect}
If $3E - 56\epsilon > 2E$, then there are at most three affected clusters, and they are adjacent in $\calG$.
\end{claim}

\begin{proof}
Suppose that $C_1,\dots, C_4$ are four clusters appearing in that order along $\lambda(a,b)$.  By item (1) of Lemma \ref{lem:cluster close}, since $d(C_i, C_{i+1})\geq E$, we have $d(s(C_i), s(C_{i+1}))> E - 14\epsilon$, for $i = 1,2,3$.

Since the shadows are disjoint segments of the geodesic $\lambda(a,b)$, we have that
$$d(s(C_1), s(C_4)) \geq d(s(C_1), s(C_2)) + d(s(C_2), s(C_3))+d(s(C_3), s(C_4)) \geq 3E -42\epsilon.$$
Hence if $3E-56\epsilon \geq 2E$, then item (1) of Lemma \ref{lem:cluster close} says that 
$$d(C_1,C_4) \geq d(s(C_1),s(C_4)) - 14\epsilon \geq 3E-56\epsilon > 2E.$$

This is impossible if $z$ is $E$-close to both $C_1,C_4$.  Hence there are at most three affected clusters and they are adjacent vertices in $\calG$.
\end{proof}



Let $C'$ denote the cluster formed by $z$ and the affected clusters from $(a,b;\calY)$.  Note that the membership of all other $(a,b;\calY)$-clusters are unchanged by the addition of $z$.  Let $C_1$ and $C_2$ be the clusters adjacent to the affected clusters in $\calG$ on the side of $a$ and $b$, respectively.  As in the split cluster case above, Lemma \ref{lem:cluster arrangement} implies that $C_1$ separates all clusters on the $a$-side of $\calG$ from $C'$, and similarly for $C_2$ and clusters on the $b$-side.  Hence the subgraph of $\calG$ connecting the cluster containing $a$ to $C_1$ is the subgraph of $\calG'$ connecting the (same) cluster containing $a$ to $C_1$, and similarly for the $b$-side of $\calG$ and $\calG'$.  As such, all of the corresponding components of $T_e$ and $T'_e$ are exactly the same, and have positive integer length by construction, verifying properties \eqref{item:integer length}, \eqref{item:stable pairs}, and \eqref{item:identical pairs} of Definition \ref{defn:stable decomp} for these components.

The forests $T_e$ and $T'_e$ have some remaining components which we need to decompose.  In $T_e$, there are the at-most $2$ components of $T_e$ which connect the affected clusters.  By assumption, these segments are at most $O(\epsilon)$ long, so we can declare them to be unstable components.

We now explain how to handle the two remaining components.  Let $A$ be the affected cluster adjacent to $C_1$ in $\calG$.  It is enough to show that the geodesic segments $\lambda(C_1,A)$ and $\lambda(C_1,C')$ are $O(\epsilon)$-Hausdorff close, since then the lengths of $\lambda_0(C_1,A)$ and $\lambda_0(C_1,C')$ are the same up to error $O(\epsilon)$, and we can declare a small end segment of each unstable, so that their (stable) complements have the same (positive integer) length and are $O(\epsilon)$-close.  But this endpoint proximity statement is again a consequence of Lemma \ref{lem:cluster close}, since the endpoints of $\lambda(C_1,A)$ and $\lambda(C_1,C')$ are within $7\epsilon$ of the endpoint of $s(C_1)$ on $\lambda(a,b)$, while $s(C')$ contains $s(A)$ and its endpoint is at most $(E+\epsilon)$-away.  Finally, the bijection $\beta:\pi_0(T-T_s) \to \pi_0(T'-T'_s)$ required for \eqref{item:adjacency} is defined similar to the split cluster case (a), with the proof of condition \eqref{item:endpoint condition} being basically the same.  For condition \eqref{item:cluster identify}, observe again the $\calY \cap (\calY \cup \{z\}) = \calY$.  Assuming that $y \in \calY$ is contained in an affected cluster in the $(\calY; \{a,b\})$ setup---since otherwise, it is contained in an identical cluster in the $(\calY \cup \{z\};\{a,b\})$ setup---then it corresponds to the single component which gets identified with the component of $T'- T'_s$ corresponding to the single cluster containing both $y$ and $z$.  This completes the proof of (1) in the affected cluster case.  See Figure \ref{fig:affected cluster} for a heuristic picture.

The argument for (2) in the affected cluster case is essentially identical to the split cluster version: All but a few components of $T_s \subset T_e$ and $T'_s \subset T'_e$ are identical and thus have identical $(r_1,r_2)$-thickenings.  And the remaining components of $T_s, T'_s$ are constructed by pairing the corresponding components of $T_e, T'_e$ adjacent to or connecting the affected clusters, and then either declaring them unstable or modifying them pairwise to achieve the length bound.  As above, this can be accomplished post-thickening; see Figure \ref{fig:affected cluster}.  This completes the proof of the theorem.

\end{proof}

\end{document}
